\section{Conclusion}
\label{sec:conclusion}

TaichiDough connects a calibrated RGB-D recording and measured two-spatula motion to floor-aware volumetric reconstruction, MLS-MPM replay, a differentiable partial-view objective, and bounded parameter optimization. On one Episode~18 sequence, the preliminary five-parameter history reduced its training objective by $61.8\%$ over 15 accepted updates and reached a late-stage plateau. A separate fixed-configuration run improved a strict visible-geometry score over frozen initial geometry during both its fitting interval and later frames of the same recorded replay. These results support the implemented observation-conditioned fitting pipeline, but not state-of-the-art identification accuracy or unique recovery of intrinsic dough properties.

The estimated values remain conditional on the camera and floor calibration, reconstructed hidden volume, registered tool geometry, contact and friction model, constitutive parameterization, numerical discretization, and selected trajectory. Future work will add synchronized force or torque sensing, repeated dough preparations, deliberately informative robot-controlled motions, explicit tool occlusion, and evaluation on independent manipulations. Once these measurements establish predictive validity beyond a single replay, the simulator can support synthetic-data generation and bimanual shaping policies.
